\documentclass[UTF8]{ctexart}

\hyphenpenalty=10000
\exhyphenpenalty=10000
\emergencystretch=1em
\usepackage{microtype}
\usepackage[T1]{fontenc}
\usepackage{lmodern} 
\usepackage{fix-cm}    
\usepackage{microtype}

\usepackage{graphicx}
\usepackage{listings}
\usepackage{xcolor}
\usepackage{booktabs}
\usepackage{array}
\usepackage{amsmath}
\usepackage{geometry}
\geometry{a4paper, margin=2.5cm}

\lstset{
  language=Python,
  basicstyle=\small\ttfamily,
  keywordstyle=\color{blue},
  commentstyle=\color{gray},
  stringstyle=\color{red},
  showstringspaces=false,
  breaklines=true,
  frame=single,
  numbers=left,
  numberstyle=\tiny\color{gray},
  tabsize=4,
  columns=flexible,
}

\begin{document}

\begin{center}
{\LARGE\textbf{人工智能实验报告}}\\[8pt]
{\Large\textbf{实验四} \quad 综合实验——基于骨骼的动作识别}
\end{center}

\vspace{8pt}

\begin{center}
\begin{tabular}{
    >{\raggedright\arraybackslash}p{4.8cm}
    >{\centering\arraybackslash}p{4.8cm}
    >{\raggedleft\arraybackslash}p{4.8cm}
}

学院：\underline{\makebox[3.6cm][c]{计算机与通信工程学院}} &
专业：\underline{\makebox[3.6cm][c]{计算机科学与技术}} &
班级：\underline{\makebox[3.6cm][c]{计 234 班}} \\
[6pt]
姓名：\underline{\makebox[3.6cm][c]{秦瑞}} &
学号：\underline{\makebox[3.6cm][c]{U202342492}} &
日期：\underline{\makebox[3.6cm][c]{2026 年 6 月 19 日}} \\
\end{tabular}
\end{center}

\vspace{12pt}

\section{实验目标}

\begin{enumerate}
  \item 数据处理能力：掌握骨骼关键点数据的清洗、归一化及特征提取方法，
        理解关节点坐标与姿态表征之间的关系，学会将原始骨骼数据转换为适用于深度学习模型的输入格式。
  \item 模型构建与训练：
        熟悉卷积神经网络（CNN）与长短期记忆网络（LSTM）的混合架构在时序动作识别任务中的应用，
        掌握 PyTorch 框架下模型搭建、训练与评估的完整流程，能根据数据特点和算力条件调整模型结构。
  \item 实验流程把控：完成从数据集加载、预处理、模型训练到测试评估的全流程操作，
        学会记录训练日志，分析训练曲线并诊断模型性能问题。
  \item 创新与优化：尝试不同的数据预处理策略、模型架构选择和超参数配置，探索提升骨骼动作识别准确率的方法。
  \item 扩展研究：通过多模态融合、实时骨骼动作识别和异常行为检测等方向的探索，
        深入了解前沿技术在实际场景中的应用潜力。
\end{enumerate}

\section{实验内容}

\subsection{任务分析}

本实验的任务是基于骨骼关键点数据进行人体动作识别。
输入数据为 CSV 格式的骨骼关节点坐标，
每个样本包含多帧、多个关节的三维坐标信息，共涵盖 8 类人体动作。

核心挑战包括：

\begin{itemize}
  \item 数据质量参差不齐：原始 CSV 据中包含无效帧和干扰受试者
        （如 \texttt{subject000}、\texttt{subject030}、\texttt{subject031} 等），需要进行过滤清洗。
  \item 时序特征提取：骨骼数据本质上是时空序列，需要同时捕捉空间上的关节配置和时间上的运动模式。
  \item 数据表示转换：需要将原始三维坐标转换为适合 CNN 输入的图像化表示（RGB 通道）。
  \item 序列长度不一致：不同样本的帧数不同，需要统一采样和裁剪到固定长度。
\end{itemize}

\subsection{实现工具}

\begin{itemize}
  \item 编程语言：Python 3
  \item 深度学习框架：PyTorch，用于模型构建、训练和推理
  \item 数据处理：NumPy、Pandas，用于 CSV 数据读取和数值计算
  \item 数据集划分：Scikit-learn（\texttt{train\_test\_split}），用于分层采样的训练/测试集划分
  \item 可视化：Matplotlib，用于绘制训练损失曲线和准确率曲线
  \item 运行环境：Jupyter Notebook 以及 Google Colab 提供的 NVIDIA Tesla T4 GPU
\end{itemize}

\subsection{实现方案}

整体方案采用 CNN-LSTM 混合架构进行骨骼动作识别，分为三个模块：

\begin{enumerate}
  \item \textbf{数据预处理模块}（\texttt{Dataset.py}）：
        读取CSV骨骼数据 $\rightarrow$ 过滤无效帧和干扰受试者 $\rightarrow$ 
        提取19个关键关节 $\rightarrow$ 以脊柱基部为中心归一化 $\rightarrow$ 
        Sigmoid变换转为RGB表示 $\rightarrow$ 时间采样和裁剪至64帧 $\rightarrow$ 
        80/20分层训练/测试划分。最终每个样本的张量形状为 $[C=3, T=64, V=19]$。
  \item \textbf{模型模块}（\texttt{Model.py}）：采用 CNN-LSTM 混合网络。
        首先通过二维卷积层提取每帧的空间特征，经自适应平均池化压缩关节维度，
        再通过双向LSTM捕捉时序依赖，最后通过全连接层输出 8 类动作分类。
  \item \textbf{训练评估模块}（\texttt{Engine.py}）：
  使用交叉熵损失函数和 Adam 优化器进行训练，每个 epoch 分别在训练集和测试集上评估性能，记录损失和准确率。
\end{enumerate}

\subsection{实现步骤}

\subsubsection{数据预处理}

数据处理流程在 \texttt{Dataset.py} 中实现，主要步骤如下：

\textbf{（1）数据加载与过滤}

从 CSV 文件读取原始骨骼数据，跳过表头。
过滤掉干扰受试者（\texttt{subject000}、\texttt{subject030}、\texttt{subject031}）的数据。
从原始特征中提取 19 个关键关节的三维坐标：

\begin{verbatim}
joint_index = [6, 7, 8, 13, 14, 15, 20, 21, 22, 27, 28, 29,
               34, 35, 36, 41, 42, 43, 48, 49, 50, 55, 56, 57,
               62, 63, 64, 174, 175, 176, 181, 182, 183, 188, 189,
               190, 195, 196, 197, 307, 308, 309, 314, 315, 316,
               321, 322, 323, 335, 336, 337, 342, 343, 344, 349,
               350, 351]
\end{verbatim}

\textbf{（2）时间采样与裁剪}

每隔 4 帧采样一次，将原始序列压缩。统计有效帧数（非零帧），
使用 \texttt{ValidCropResize} 函数以 $0.95$ 的裁剪比例进行裁剪，
并通过双线性插值统一 resize 到 64 帧。

\textbf{（3）归一化与 RGB 转换}

以脊柱基部（spine base）关节点为原点，将所有关节坐标进行中心化归一化：
\[
\text{data\_centered} = \text{data} - \text{spinebase}
\]

然后通过类 Sigmoid 变换将坐标映射到 RGB 图像空间：
\[
\text{RGB} = \left\lceil \frac{255}{1 + e^{-0.01 \times \text{data\_centered}}} \right\rceil
\]

\textbf{（4）数据集划分}

使用 \texttt{train\_test\_split} 以 $80/20$ 比例进行分层划分（stratified split），
保证各类别在训练集和测试集中的比例一致。最终每个样本的形状为 $[C=3, T=64, V=19]$，共 8 个动作类别。

\subsubsection{模型构建}

CNN-LSTM 混合模型在 \texttt{Model.py} 中定义，网络结构如下：

\begin{enumerate}
  \item \textbf{卷积层}：
        \begin{center}
          \texttt{Conv2d(3, 64, kernel\_size=3, padding=1) + ReLU}
        \end{center}
        激活。输入通道为 3（RGB），输出64个特征图，通过 $3 \times 3$ 卷积核提取每帧的空间特征。
  \item \textbf{自适应池化层}：
        \begin{center}
          \texttt{AdaptiveAvgPool2d((64, 1))}
        \end{center}
        将 19 个关节维度压缩为 1，得到形状 $[B, 64, 64, 1]$。
  \item \textbf{LSTM层}：
        \begin{center}
          \texttt{LSTM(input\_size=64, hidden\_size=128, num\_layers=2, bidirectional=True)}
        \end{center}
        双向 LSTM 捕捉前后向时序依赖，输出维度为 $128 \times 2 = 256$。
  \item \textbf{全连接层}：
        \begin{center}
          \texttt{Linear(256, 64)} + ReLU + \texttt{Linear(64, 8)}
        \end{center}
        将 LSTM 的最终时间步输出映射到8类动作分类。
\end{enumerate}

模型前向传播过程：
\[
x \xrightarrow{\text{Conv2D + ReLU}}
  \xrightarrow{\text{AdaptiveAvgPool}}
  \xrightarrow{\text{Permute}}
  \xrightarrow{\text{BiLSTM}}
  \xrightarrow{\text{FC + ReLU}}
  \xrightarrow{\text{FC}}
\hat{y}
\]

此外，\texttt{Model.py} 中还定义了 \texttt{ResNetBaseline} 模型作为对比基准，
使用预定义的 ResNet-18 结构替换最后的全连接层。

\subsubsection{训练与评估}

训练流程在 \texttt{Engine.py} 中实现：

\begin{itemize}
  \item \textbf{损失函数}：交叉熵损失（CrossEntropyLoss），适用于多分类任务。
  \item \textbf{优化器}：Adam 优化器，学习率 \texttt{lr=0.0007}，权重衰减 \texttt{weight\_decay=4e-6}。
  \item \textbf{训练配置}：批大小（batch size）为 128，训练 300 个 epoch。
  \item \textbf{评估指标}：训练损失和分类准确率。
\end{itemize}

每个epoch的训练过程：
\begin{enumerate}
  \item 模型切换为训练模式，遍历训练数据，前向传播计算损失，反向传播更新权重。
  \item 模型切换为评估模式，在测试集上计算准确率。
  \item 记录当前 epoch 的训练损失和测试准确率。
\end{enumerate}

\subsubsection{实验结果}

实验结果如图所示：

\begin{center}
  \includegraphics[width=0.9\textwidth]{Build/img_000.png}
\end{center}

从训练曲线可以观察到：

\begin{itemize}
  \item \textbf{训练损失}：在训练初期迅速下降，随后逐渐趋于平稳，表明模型学习到了有效的特征表示。
  \item \textbf{测试准确率}：最佳测试准确率达到 $94.14\%$，最终测试准确率为 $86.08\%$。
        后期准确率出现波动属于正常现象，
        可能由学习率偏大导致模型在最优点附近震荡，或过拟合后泛化能力下降所致。
\end{itemize}

\section{实验总结}

\textbf{（1）模型性能分析}：CNN-LSTM 混合模型在骨骼动作识别任务上取得了较好的效果。
CNN 层负责提取每帧的空间特征（关节间的空间关系），LSTM 层负责捕捉帧间的时序动态（动作的时间演化过程），
两者结合能够有效建模骨骼数据的时空特性。

\textbf{（2）模型参数影响}：
\begin{itemize}
  \item 学习率 \texttt{lr=0.0007} 的选择对训练稳定性至关重要，过大会导致震荡，过小会收敛缓慢。
  \item 权重衰减 \texttt{weight\_decay=4e-6} 提供了正则化效果，一定程度上抑制了过拟合。
  \item 双向 LSTM 相比单向 LSTM 能够利用未来帧的信息，提升了分类性能。
  \item 2 层 LSTM 的堆叠增加了模型的表达能力，但也增加了计算开销。
\end{itemize}

\textbf{（3）存在的问题}：
\begin{itemize}
  \item 训练后期准确率出现波动（从 $94.14\%$ 降至 $86.08\%$），
        可能需要引入学习率衰减策略（如 CosineAnnealing 或 ReduceLROnPlateau）来稳定训练。
  \item 数据预处理中简单的 Sigmoid 变换可能丢失部分精细的空间信息，可探索更高级的骨骼表示方法。
  \item 当前模型未使用数据增强（如随机旋转、缩放、时间扰动等），有一定的性能提升空间。
\end{itemize}

\textbf{（4）改进建议}：可以尝试引入学习率调度器、数据增强策略、多尺度特征融合或注意力机制（如 ST-GCN）来进一步提升识别性能。
同时，ResNetBaseline的对比实验可以作为消融研究的一部分，验证CNN-LSTM架构的有效性。

\section{扩展研究}

\textbf{基于骨骼的动作识别在实际场景中的应用——实时姿态估计与行为分析}

基于骨骼的动作识别技术在实际应用中具有广泛前景：

\textbf{（1）多模态融合}：将骨骼数据与 RGB 视频、深度图像等多模态信息融合，可以显著提升识别准确率。
例如，OpenPose 或 MediaPipe 等工具可以实时提取人体骨骼关键点，结合视频外观特征进行联合建模。

\textbf{（2）轻量化与实时部署}：对于实时应用场景（如安防监控、人机交互），需要对模型进行轻量化处理。
可以通过知识蒸馏、模型剪枝或量化技术，将 CNN-LSTM 模型部署到边缘设备上，实现实时推理。

\textbf{（3）异常行为检测}：在安防、医疗等场景中，基于骨骼的动作识别可用于异常行为（如跌倒、暴力行为）的自动检测。
这需要在正常行为数据的基础上，引入异常检测算法（如 One-Class SVM、自编码器等）。

\textbf{（4）图卷积网络（GCN）}：将人体骨骼建模为图结构，使用图卷积网络（如 ST-GCN）可以直接在骨骼拓扑上进行特征学习，
相比 CNN-LSTM 方法更能捕捉关节间的拓扑关系，是当前骨骼动作识别领域的主流方法。

\end{document}
